\section{Calibrating the Pair on Real Data}\label{sec:calib}

A pair $(m_0,m_1)$ is a statement about the sensor, the window length and the resolution: the hit rate that structure reaches and the hit rate that clutter does not exceed. Section~\ref{sec:seq} makes the user state it. This section asks whether we can state it from data. We set the procedure down in a pre-registration before the first calibration measurement and applied it to calibration sequences that are separate from the test sequences of Section~\ref{sec:realeval}. The pre-registration was amended after that measurement and after an independent review of it, before the measurement that the tables below report and before any test-sequence file was opened. The first amendment changed the time base of the LILocBench poses, the use of the semantic classes of KITTI 00, the survivorship of cells with few touched windows, the bound on the skew between the two lasers of the pose check and the wall-thickness gate. The second amendment set the quality gates, the rounding direction, the one-second window, the 3D miss branch and the move of Keble~02 to the calibration set. The procedure below is the amended one. We report the outcome. \emph{The outcome is that nothing froze.} The shipped pair is the pair that the procedure would give if the failed quality gate were waived, and it is labelled provisional wherever it appears.

\subsection{Procedure}\label{subsec:calibproc}

\noindent\textbf{Data and split.} The calibration set holds the LILocBench \citep{trekel2025lilocbench} sequences \texttt{mapping}, \texttt{static\_0} and \texttt{dynamics\_0} (2D, a wheeled robot in an office), the KTH benchmark \citep{zhang2023dynamic} sequence KITTI 00 (3D, 64 beams, a single drive) and, after a re-split described below, the Oxford Spires \citep{tao2024spires} sequence Keble~02 (3D, 64 beams, one site). The test set holds the sequences of Section~\ref{sec:realeval}. The split, the statistic, the thresholds and the rule below, in their amended form, were committed before any test-sequence file was opened; the files of the test set were recorded with size and modification time before and after each phase, and every later read of a test file is listed in an access log. The re-split moved Keble~02, the lowest-numbered Keble sequence that has a ground-truth trajectory (a rule that does not look at content), from the test set to the calibration set, in an amendment committed before any test file was opened. It was needed because KITTI 00 alone, a short drive, cannot show whether most structure voxels can be decided by a sequential test.

\noindent\textbf{Cells and labels.} For each cell touched in at least five windows, the hit rate is $r=H/N$ with $H$ the windows with a hit and $N$ the touched windows. A miss is a window in which a ray passed through the cell to a return beyond it; a beam without a return adds nothing. Structure cells are those that the reference holds as occupied (LILocBench: occupied in \texttt{static\_0}; KITTI 00: static voxels above the ground; Keble~02: voxels of the terrestrial-laser-scanner map, restricted to the area around the trajectory). Clutter cells are those in which something moves (LILocBench: cells hit in \texttt{dynamics\_0} that were free in the mapping reference; KITTI 00: voxels that hold a point flagged dynamic; Keble~02: hits that the terrestrial map does not contain). The window is one second, the length of the windows of the node: 15 scans at 15~Hz in 2D, 10 scans at 10~Hz in 3D. A first calibration run used ten scans per window (0.67~s at 15~Hz); the one-second condition was added before the freeze step.

\noindent\textbf{Rule.} Let $q_s$ be the 10th percentile of the hit rate over structure cells and $q_c$ the 90th percentile over clutter cells. The pre-registered escalation rule has three steps. (i) If $q_s-q_c\ge0.15$, take $m_0=\lceil q_c\rceil_{0.05}$ and $m_1=\lfloor q_s\rfloor_{0.05}$, rounded toward each other so that the stated error rates cover the measured percentiles, with $\alpha=\beta=0.01$. (ii) Otherwise measure the gap again with a refined miss that does not count a miss next to the endpoint of a grazing beam (3D only). (iii) Otherwise add a second statistic on run lengths (Option C). In every branch the remaining overlap is reported as irreducible. The truncation $n_{\max}$ is the smallest multiple of 16 for which at most 3\% of cells at the zero-drift rate $m^\ast$ are still undecided under the exact rule.

\noindent\textbf{Quality gates.} A hit-rate percentile is only as good as the map that carries it. Two gates were fixed from sensor specifications in the second amendment, after the first calibration measurement (whose gate looked at the extent of a patch only) and before the measurement that the tables report. G1 measures the thickness $\sqrt{12\lambda_{\min}}$ of linear (2D) or planar and vertical (3D) patches of the accumulated points. It passes if the median and the 90th percentile of the thickness stay under limits derived from the error of the sensor and of the poses (Table~\ref{tab:gates}); the limit of the 90th percentile is twice that of the median. G2 measures pose consistency: each return is compared with the returns of scans one to four seconds away, and the median and 90th percentile of the distance to the nearest returned point must stay within one and two cells of the finest grid of the sensor class. A sequence that fails a gate is excluded from the headline tables; the 2D pair needs all three LILocBench sequences to pass.

\subsection{Result}\label{subsec:calibresult}

% AUTO-GENERATED by tools/make_v2_artifacts.py from results_p1b/p1b_gates.csv -- do not edit
\begin{table}[t]
\centering
\caption{Quality gates of the calibration sequences (P1b, pre-registered thresholds from the sensor specifications). G1: wall thickness of planar patches (median and 90th percentile of the thickness, metres); G2: leave-one-out pose consistency (median and 90th percentile of the distance, metres). Limits are written under the column names. A sequence that fails a gate is excluded from the headline tables. \emph{No sequence passes}, so the pair derived from them is not frozen.}
\label{tab:gates}
\footnotesize
\setlength{\tabcolsep}{4pt}
\begin{tabular}{llcccccc}
\toprule
Sequence & Class & G1 patches & G1 median / p90 & G1 limits & G2 median / p90 & G2 limits & G1, G2 \\
\midrule
\texttt{lil\_mapping} & 2D & 590 & 0.1930 / 0.2506 & 0.1 / 0.2 & 0.0012 / 0.0048 & 0.05 / 0.1 & FAIL, PASS \\
\texttt{lil\_static\_0} & 2D & 391 & 0.1705 / 0.2545 & 0.1 / 0.2 & 0.0013 / 0.0060 & 0.05 / 0.1 & FAIL, PASS \\
\texttt{lil\_dynamics\_0} & 2D & 470 & 0.1832 / 0.2285 & 0.1 / 0.2 & 0.0021 / 0.0276 & 0.05 / 0.1 & FAIL, PASS \\
\texttt{kitti00} & 3D & 41 & 0.3720 / 0.4713 & 0.2 / 0.4 & 0.0387 / 0.1113 & 0.1 / 0.2 & FAIL, PASS \\
\texttt{keble02} & 3D & 2251 & 0.3212 / 0.4627 & 0.2 / 0.4 & 0.0174 / 0.0415 & 0.1 / 0.2 & FAIL, PASS \\
\bottomrule
\end{tabular}
\end{table}


\noindent\textbf{The gates.} G2 passes on all \CalGateSeq{} calibration sequences. G1 fails on all \CalGateGoneFail{}: the median patch thickness on the three 2D sequences is between \CalGoneMedLo{} and \CalGoneMedHi{}~m against a limit of \CalGoneLimMed{}~m (Table~\ref{tab:gates}). We cannot say whether the cause is the poses, the time base of the scans or the walls themselves. A search for the time offset between scans and poses was run on the \texttt{mapping} sequence by a procedure written before use; a first calibration run without it measured pose-timing smear more than hit rates. The gate also fails on an exploratory run without the offset. Because no sequence passes, every table below is printed for the record and is not a headline table.

% AUTO-GENERATED by tools/make_v2_artifacts.py from results_p1b/p1b_gap.csv, results_p1/p1_gap.csv -- do not edit
\begin{table}[t]
\centering
\caption{Hit-rate gap of structure against clutter on the calibration sequences. $q_s$: 10th percentile of the per-cell hit rate over structure cells; $q_c$: 90th percentile over clutter cells; both for cells touched in at least five windows. The pre-registered rule needs a gap $q_s-q_c\ge0.15$. Block intervals over 10~m tiles in brackets. \emph{LIL}: LILocBench, 2D, structure of \texttt{static\_0} against hits in free space of \texttt{dynamics\_0}. \emph{KITTI 00}: KTH benchmark, 3D, structure voxels against voxels that hold a moving point (first calibration run, ten scans per window). \emph{Keble 02}: Oxford Spires, 3D, structure against TLS-absent hits. $m_0,m_1$: the rounding rule of Section~\ref{sec:calib}, defined only when the gap passes; all rows fail a quality gate (Table~\ref{tab:gates}).}
\label{tab:gap}
\footnotesize
\setlength{\tabcolsep}{3.5pt}
\resizebox{\linewidth}{!}{%
\begin{tabular}{llcccccc}
\toprule
Data & Cell, window & $q_s$ & $q_c$ & Gap [block interval] & $\ge0.15$ & $(m_0,m_1)$ & Cells (struct. / clutter) \\
\midrule
LIL & 0.05 m, 15 scans & 0.7130 & 0.1725 & 0.5404 [0.4435, 0.5795] & yes & (0.20, 0.70) & 10{,}126 / 23{,}888 \\
LIL & 0.10 m, 15 scans & 0.7222 & 0.2381 & 0.4841 [0.3746, 0.5183] & yes & (0.25, 0.70) & 4{,}582 / 6{,}220 \\
KITTI 00 & 0.2 m, 10 scans & 0.1250 & 0.1667 & -0.0417 [-0.1071, -0.0238] & no & -- & 121{,}651 / 5{,}791 \\
Keble 02 & 0.2 m, 10 scans, today's miss & 0.4667 & 1.0000 & -0.5333 [-0.6000, -0.4762] & no & -- & 195{,}641 / 15{,}864 \\
Keble 02 & 0.2 m, 10 scans, refined miss & 0.5357 & 1.0000 & -0.4643 [-0.5333, -0.4000] & no & -- & 192{,}255 / 13{,}990 \\
\bottomrule
\end{tabular}}
\end{table}


\noindent\textbf{The hit-rate gap.} In 2D the gap is large: $q_s=\CalQs$, $q_c=\CalQc$, a gap of \CalGap{} with a block-bootstrap interval of [\CalGapLo, \CalGapHi] over \CalNstruct{} structure and \CalNclutter{} clutter cells (Table~\ref{tab:gap}). The rounding rule gives $(m_0,m_1)=(\CalMzero,\CalMone)$. The intervals of $q_s$ and $q_c$ cross rounding boundaries, so the pair at both ends of the interval is $(\CalLowEndZero,\CalLowEndOne)$ and $(\CalHighEndZero,\CalHighEndOne)$; both are evaluated in Section~\ref{sec:realeval}. At the coarser resolution of 0.1~m the rule gives $(\CalMzeroCoarse,\CalMoneCoarse)$. With the pair $(\CalMzero,\CalMone,0.01,0.01)$ the derived quantities are those of Table~\ref{tab:sprtdefaults}, among them $n_{\max}=\CalNmaxCond$ and an always-hit latency of \CalLatCond{} windows. In 3D the gap is negative under both miss definitions: $q_s=\CalKebQs$ against $q_c=\CalKebQc$ on Keble~02, a gap of \CalKebGapToday{} with today's miss and \CalKebGapRefined{} with the refined miss, and $\CalKthGap$ on KITTI 00 ($q_s=\CalKthQs$, $q_c=\CalKthQc$). The Keble~02 result is not evidence that hit rates do not separate people from structure in 3D: \CalKebAbove\,\% of the clutter cells lie more than half a metre above the device, so the clutter class is mostly not walking people. On KITTI 00 a structure voxel is touched in a median of \CalKittiMed{} windows (\CalKittiWindows{} windows close), too few for a sequential test to decide most of them; a lifelong mapper needs revisits.

% AUTO-GENERATED by tools/make_v2_artifacts.py from results_p1b/p1b_matched_band.csv, results_p1/p1_matched_band.csv -- do not edit
\begin{table}[t]
\centering
\caption{Evidence for a second statistic (Option C) on the calibration data: cells matched on hit rate (bands of width 0.2), lag-1 autocorrelation of the window outcomes. AUC: probability that a structure cell has a larger lag-1 autocorrelation than a clutter cell; above 0.5 means structure has the longer runs. Pre-registered reading: support for Option C if the AUC is at least 0.65 or at most 0.35 in a band with at least 100 cells in each group.}
\label{tab:optc}
\footnotesize
\setlength{\tabcolsep}{4pt}
\resizebox{\linewidth}{!}{%
\begin{tabular}{llcccccl}
\toprule
Data & Hit-rate band & Struct. cells & Clutter cells & Lag-1, struct. & Lag-1, clutter & AUC & Reading \\
\midrule
LIL, 0.05 m, 15 scans & $[0.2,0.4)$ & 99 & 1{,}382 & 0.329 & 0.360 & 0.493 & fewer than 100 cells \\
 & $[0.4,0.6)$ & 337 & 172 & 0.474 & 0.274 & 0.665 & support, structure longer \\
 & $[0.6,0.8)$ & 1{,}288 & 91 & 0.468 & 0.231 & 0.720 & fewer than 100 cells \\
Keble 02, 0.2 m, 10 scans & $[0.2,0.4)$ & 7{,}782 & 2{,}055 & 0.311 & 0.010 & 0.736 & support, structure longer \\
 & $[0.4,0.6)$ & 13{,}988 & 2{,}102 & 0.396 & 0.020 & 0.771 & support, structure longer \\
 & $[0.6,0.8)$ & 27{,}402 & 2{,}526 & 0.368 & -0.006 & 0.776 & support, structure longer \\
KITTI 00, 0.2 m, 10 scans & $[0.2,0.4)$ & 25{,}052 & 137 & 0.154 & 0.360 & 0.338 & support, structure shorter \\
 & $[0.4,0.6)$ & 22{,}018 & 134 & 0.260 & 0.339 & 0.468 & between \\
 & $[0.6,0.8)$ & 19{,}036 & 147 & 0.179 & 0.272 & 0.436 & between \\
\bottomrule
\end{tabular}}
\end{table}


\noindent\textbf{A second statistic.} The pre-registered reading counts support for Option C if the lag-1 autocorrelation of the window outcomes separates structure from clutter at matched hit rate, with an AUC of at least 0.65 or at most 0.35 in a band with at least 100 cells in each group (Table~\ref{tab:optc}). It does, with opposite signs: in LILocBench (AUC \CalOptcLil) and Keble~02 (\CalOptcKebLo{} to \CalOptcKebHi) structure has the longer runs; in KITTI 00 (\CalOptcKth) it has the shorter ones. A statistic whose sign changes between datasets is not a fixed rule, and we did not implement Option C. This is a measured outcome, not a decision that the gap rule made.

\noindent\textbf{The pair on its own calibration data.} Replayed through the engine on the three one-second 2D sequences, the shipped pair with the Periodic veto on leaves \CalRepLe\,\% of the \CalRepNle{} clutter cells whose observed hit rate is at most $m_0$ ever Static, \CalRepGt\,\% of those above $m_0$, and \CalRepAll\,\% of all \CalRepN{} clutter cells. These per-cell rates exceed the per-history level $\alpha=0.01$ for the reasons of Section~\ref{subsec:promise}; they are the calibration rows that the held-out rates of Section~\ref{sec:realeval} are read beside.

\subsection{What froze and what did not}\label{subsec:froze}

Nothing froze. The 2D pair $(\CalMzero,\CalMone,0.01,0.01)$ is the result of the rule with the failed gate waived, not a freeze. The 3D rule gave no pair, and the 3D backend falls back on the 2D pair, which is \emph{not calibrated for 3D}. The shipped defaults are therefore provisional: the node logs this at start and publishes it on its diagnostics topic, and a test of the library fails if the library defaults, the shipped parameter files and the pre-registration record disagree. The pair applies only to one-second windows, the resolutions measured and the sensors measured. The shipped parameter files set \code{window\_mode} to \code{time} with one-second windows. The code defaults of the node are \code{window\_mode} \code{scans} with \code{layer\_interval} $10$, which are not one-second windows, so the pair holds for the shipped parameter files and not for a node started on the code defaults. Re-deriving the pair for another sensor needs a structure reference and per-cell clutter labels, one-second windows, the hit-rate probe (\code{hitrate\_probe}) and the gates G1 and G2. The gate G1 failed here even on benchmark data with reference poses, and without labels $(m_0,m_1)$ can only be chosen by judgement. We evaluate the point pair and both ends of its interval in Section~\ref{sec:realeval} so that the reader sees how much the answers depend on the pair.
